Aunque el lenguaje Java abstrae la gestión directa de la memoria física delegando esta tarea al
a su sistema automatizado, el desarrollo de este sistema implementa un control explícito sobre las estructuras de datos y los
recursos del Sistema Operativo para evitar fugas de memoria (\textit{memory leaks}) y la saturación del espacio.

A continuación, se describen las dos estrategias principales utilizadas para la optimización de la memoria y los recursos.

\subsubsection{Liberación de recursos de Entrada/Salida}
Al interactuar con el sistema de archivos del Sistema Operativo, cada archivo abierto consume un Descriptor de Archivo
(\textit{File Descriptor}) en la tabla de procesos del kernel. Si estos descriptores y los búferes
de memoria asociados no se liberan, el sistema puede colapsar. 

Para prevenir esto, el proyecto hace uso del paradigma \textit{try-with-resources},
el cual garantiza que las conexiones de memoria y los descriptores del SO se cierren atómicamente
tan pronto como el bloque termine su ejecución, incluso si ocurre una interrupción de software (excepción).

\vspace{-0.5cm}
\begin{figure}[H]
    \centering
    \caption{Gestión segura de memoria y descriptores en I/O}
    \label{code:gestion_memoria_io}
    \begin{lstlisting}[language=Java]
private List<Alumno> cargarAlumnosDesdePDF() throws Exception {
    // ...
    // El DirectoryStream asegura la liberacion del iterador de directorio en el SO
    try (DirectoryStream<Path> stream = Files.newDirectoryStream(CARPETA_PDF, "*.pdf")) {
        for (Path pdf : stream) {
            System.out.println("Procesando: " + pdf.getFileName());

            try {
                // El bloque garantiza que el buffer de memoria del PDDocument
                // sea destruido inmediatamente tras extraer el texto.
                try (PDDocument doc = PDDocument.load(pdf.toFile())) {
                    PDFHorarioStripper stripper = new PDFHorarioStripper();
                    stripper.getText(doc);
                    alumnos.add(new Alumno(stripper.getNombreAlumno(), stripper.getHorario()));
                }
            } catch (IOException e) {
                // ... manejo de excepciones
            }
        }
    }
    return alumnos;
}
    \end{lstlisting}
\end{figure}

Como se observa en la figura \ref{code:gestion_memoria_io}, tanto el flujo del directorio (\texttt{DirectoryStream})
como el buffer del documento en memoria RAM (\texttt{PDDocument}) están encapsulados.
Tan pronto como el recolector extrae los horarios, esa memoria es marcada como libre para el sistema.

\subsubsection{Optimización mediante EnumMap}
Para el almacenamiento de los horarios procesados, el sistema evita el uso de mapas hash tradicionales
(\texttt{HashMap}) y en su lugar implementa \texttt{EnumMap}.

\vspace{-0.5cm}
\begin{figure}[H]
    \centering
    \caption{Asignación de memoria contigua para estructuras de horarios}
    \label{code:gestion_memoria_enummap}
    \begin{lstlisting}[language=Java]
public static Map<DayOfWeek, List<Intervalo>> unirHorariosOcupados(List<Alumno> alumnos) {

    // Inicializacion optimizada en memoria
    Map<DayOfWeek, List<Intervalo>> ocupados = new EnumMap<>(DayOfWeek.class);

    for (DayOfWeek d : DayOfWeek.values()) {
        ocupados.put(d, new ArrayList<>());
    }
    // ... union de intervalos
}
    \end{lstlisting}
\end{figure}

A nivel de arquitectura, \texttt{EnumMap} se representa internamente como un arreglo plano (\textit{array})
en lugar de una estructura de nodos enlazados. Esta decisión de diseño asigna un bloque de
memoria contiguo, lo que mejora drásticamente la Localidad de Referencia Espacial en la memoria Caché del
procesador al iterar sobre los días de la semana, reduciendo los fallos de página y acelerando el cálculo de intervalos libres.